\documentclass[11pt,a4paper]{article}

% =========================================================================
% PAKET-PAKET UTAMA
% =========================================================================
\usepackage[utf8]{inputenc}
\usepackage[T1]{fontenc}
\usepackage{lmodern}
\usepackage[indonesian]{babel}
\usepackage{amsmath,amssymb,amsfonts,amsthm}
\usepackage{bm}              % Notasi tebal untuk vektor & matriks
\usepackage{geometry}
\geometry{top=2.0cm, bottom=2.0cm, left=2.0cm, right=2.0cm}
\usepackage{graphicx}
\usepackage{booktabs}        % Tabel profesional
\usepackage{tabularx}
\usepackage{microtype}
\usepackage{xcolor}
\usepackage{mdframed}        % Kotak highlight untuk definisi
\usepackage{caption}
\usepackage{subcaption}
\usepackage{float}
\usepackage{hyperref}
\hypersetup{
    colorlinks=true,
    linkcolor=blue!70!black,
    citecolor=blue!70!black,
    urlcolor=blue!70!black,
    pdfencoding=auto,
    psdextra
}
\pdfstringdefDisableCommands{%
  \def\boldsymbol#1{#1}%
  \def\bm#1{#1}%
  \def\mathbf#1{#1}%
}

% Makro sitasi fleksibel
\newcommand{\citep}[2][]{(Rasmussen \& Williams, 2006\if\relax\detokenize{#1}\relax\else, #1\fi)}
\newcommand{\citet}[2][]{Rasmussen \& Williams (2006\if\relax\detokenize{#1}\relax\else, #1\fi)}

% Lingkungan Teorema & Definisi
\theoremstyle{plain}
\newtheorem{teorema}{Teorema}[section]
\newtheorem{claim}{Claim}[section]

\theoremstyle{definition}
\newtheorem{definisi}{Definisi}[section]

\theoremstyle{remark}
\newtheorem{catatan}{Catatan}[section]

% Definisi warna kustom untuk kotak
\definecolor{dgpblue}{RGB}{26, 82, 118}
\definecolor{dgpred}{RGB}{148, 49, 38}

\newmdenv[
    linecolor=dgpblue,
    outerlinewidth=1.2pt,
    roundcorner=4pt,
    innertopmargin=6pt,
    innerbottommargin=6pt,
    innerrightmargin=10pt,
    innerleftmargin=10pt,
    backgroundcolor=blue!2,
    nobreak=true
]{mathbox}

\newmdenv[
    linecolor=dgpred,
    outerlinewidth=1.2pt,
    roundcorner=4pt,
    innertopmargin=6pt,
    innerbottommargin=6pt,
    innerrightmargin=10pt,
    innerleftmargin=10pt,
    backgroundcolor=red!2,
    nobreak=true
]{alertbox}

% Format nama pembuktian
\addto\captionsindonesian{\renewcommand{\proofname}{Proof}}
\renewcommand{\proofname}{Proof}

% Path direktori gambar
\graphicspath{{figures/}}

% =========================================================================
% INFORMASI DOKUMEN
% =========================================================================
\title{
    \vspace{-1.6cm}
    \textbf{\Large BIMBINGAN TUGAS AKHIR \#8} \\[0.2em]
}
\author{}
\date{\vspace{-0.6cm} Selasa, 3 November 2026}

\begin{document}

\maketitle
\vspace{-1.1cm}
\setlength{\abovedisplayskip}{3pt}
\setlength{\belowdisplayskip}{3pt}
\setlength{\abovedisplayshortskip}{1pt}
\setlength{\belowdisplayshortskip}{1pt}

% =========================================================================
% 1. DEEP NEURAL NETWORK
% =========================================================================
\section{Deep Neural Network}
\label{sec:dnn}

\noindent Diberikan dataset $\mathcal{D} = \{\mathbf{x}_i\}_{i=1}^N$ dengan $N$ data sampel. Setiap data ke-$i$ dinyatakan sebagai vektor baris $\mathbf{x}_i \in \mathbb{R}^{1 \times D}$ dengan $D$ fitur:
\begin{equation}
    \mathbf{x}_i = \begin{bmatrix} x_{i, 1} & x_{i, 2} & \dots & x_{i, D} \end{bmatrix} \in \mathbb{R}^{1 \times D} \label{eq:vektor_fitur_data_i}
\end{equation}
Untuk arsitektur dengan $L$ layer ($\ell = 1, \dots, L$) dan neuron $n_\ell \in \mathbb{N}^+$:
\begin{equation}
    \mathbf{x}^{(0)}_i = \mathbf{x}_i \in \mathbb{R}^{1 \times n_0}, \quad n_0 = D, \quad x^{(0)}_{i, k} = x_{i, k}
\end{equation}
di mana $\mathbf{x}^{(\ell-1)}_i \in \mathbb{R}^{1 \times n_{\ell-1}}$ ditransformasikan menjadi $\mathbf{x}^{(\ell)}_i \in \mathbb{R}^{1 \times n_\ell}$ hingga mencapai output layer $\mathbf{x}^{(L)}_i \in \mathbb{R}^{1 \times n_L}$.

% -------------------------------------------------------------------------
\subsection{Tinjau Satu Neuron}
\label{subsec:neuron_skalar}

Untuk neuron ke-$j$ pada layer $\ell$ ($j = 1, \dots, n_\ell$), komputasi terhadap data ke-$i$ dirumuskan sebagai:
\begin{align}
    z^{(\ell)}_{i, j} &= \sum_{k=1}^{n_{\ell-1}} x^{(\ell-1)}_{i, k} w^{(\ell)}_{kj} + b^{(\ell)}_j \label{eq:scalar_z} \\
    x^{(\ell)}_{i, j} &= \sigma^{(\ell)}\left( z^{(\ell)}_{i, j} \right) \label{eq:scalar_x}
\end{align}
dengan indeks dan parameter:
\begin{itemize}\setlength{\itemsep}{-2pt}\setlength{\parsep}{0pt}\setlength{\topsep}{0pt}
    \item $i \in \{1, \dots, N\}$: indeks data ke-$i$.
    \item $k \in \{1, \dots, n_{\ell-1}\}$: indeks fitur input (layer $\ell-1$).
    \item $j \in \{1, \dots, n_\ell\}$: indeks neuron / fitur output (layer $\ell$).
    \item $w^{(\ell)}_{kj} \in \mathbb{R}$: bobot fitur $k$ ke neuron $j$.
    \item $b^{(\ell)}_j \in \mathbb{R}$: bias neuron $j$.
    \item $\sigma^{(\ell)}: \mathbb{R} \to \mathbb{R}$: fungsi aktivasi element-wise.
\end{itemize}

% -------------------------------------------------------------------------
\subsection{Tinjau Satu Data}
\label{subsec:vektorisasi_layer}

Didefinisikan vektor baris data fitur ke-$i$, matriks bobot $W^{(\ell)}$, dan vektor bias untuk layer $\ell$:
\begin{align}
    \mathbf{x}^{(\ell-1)}_i &= \begin{bmatrix} x^{(\ell-1)}_{i, 1} & x^{(\ell-1)}_{i, 2} & \dots & x^{(\ell-1)}_{i, n_{\ell-1}} \end{bmatrix} \in \mathbb{R}^{1 \times n_{\ell-1}} \label{eq:vector_x_prev} \\
    W^{(\ell)} &= \begin{bmatrix}
        w^{(\ell)}_{11} & w^{(\ell)}_{12} & \dots & w^{(\ell)}_{1 n_\ell} \\
        w^{(\ell)}_{21} & w^{(\ell)}_{22} & \dots & w^{(\ell)}_{2 n_\ell} \\
        \vdots & \vdots & \ddots & \vdots \\
        w^{(\ell)}_{n_{\ell-1} 1} & w^{(\ell)}_{n_{\ell-1} 2} & \dots & w^{(\ell)}_{n_{\ell-1} n_\ell}
    \end{bmatrix} \in \mathbb{R}^{n_{\ell-1} \times n_\ell} \label{eq:matrix_W} \\
    \mathbf{b}^{(\ell)} &= \begin{bmatrix} b^{(\ell)}_1 & \dots & b^{(\ell)}_{n_\ell} \end{bmatrix} \in \mathbb{R}^{1 \times n_\ell}, \quad
    \mathbf{z}^{(\ell)}_i = \begin{bmatrix} z^{(\ell)}_{i, 1} & \dots & z^{(\ell)}_{i, n_\ell} \end{bmatrix} \in \mathbb{R}^{1 \times n_\ell} \label{eq:vectors_bz}
\end{align}
Transformasi vektor baris pada layer $\ell$ untuk data ke-$i$:
\begin{align}
    \underbrace{\mathbf{z}^{(\ell)}_i}_{(1 \times n_\ell)} &= \underbrace{\mathbf{x}^{(\ell-1)}_i}_{(1 \times n_{\ell-1})} \times \underbrace{W^{(\ell)}}_{(n_{\ell-1} \times n_\ell)} + \underbrace{\mathbf{b}^{(\ell)}}_{(1 \times n_\ell)} \label{eq:vector_z} \\
    \underbrace{\mathbf{x}^{(\ell)}_i}_{(1 \times n_\ell)} &= \sigma^{(\ell)}\left( \mathbf{z}^{(\ell)}_i \right) \label{eq:vector_x}
\end{align}

% -------------------------------------------------------------------------
\subsection{Tinjau \texorpdfstring{$N$}{N} Data}
\label{subsec:matriks_batch}

Diberikan matriks data input $X^{(\ell-1)} \in \mathbb{R}^{N \times n_{\ell-1}}$ di mana setiap baris ke-$i$ memuat vektor baris fitur dari data ke-$i$, $\mathbf{x}^{(\ell-1)}_i$:
\begin{equation}
    X^{(\ell-1)} = \begin{bmatrix}
        \mathbf{x}^{(\ell-1)}_1 \\
        \mathbf{x}^{(\ell-1)}_2 \\
        \vdots \\
        \mathbf{x}^{(\ell-1)}_i \\
        \vdots \\
        \mathbf{x}^{(\ell-1)}_N
    \end{bmatrix} = \begin{bmatrix}
        x^{(\ell-1)}_{1, 1} & x^{(\ell-1)}_{1, 2} & \dots & x^{(\ell-1)}_{1, n_{\ell-1}} \\
        x^{(\ell-1)}_{2, 1} & x^{(\ell-1)}_{2, 2} & \dots & x^{(\ell-1)}_{2, n_{\ell-1}} \\
        \vdots & \vdots & \ddots & \vdots \\
        x^{(\ell-1)}_{i, 1} & x^{(\ell-1)}_{i, 2} & \dots & x^{(\ell-1)}_{i, n_{\ell-1}} \\
        \vdots & \vdots & \ddots & \vdots \\
        x^{(\ell-1)}_{N, 1} & x^{(\ell-1)}_{N, 2} & \dots & x^{(\ell-1)}_{N, n_{\ell-1}}
    \end{bmatrix} \in \mathbb{R}^{N \times n_{\ell-1}} \label{eq:matrix_batch_X}
\end{equation}
dengan $X^{(0)} = X \in \mathbb{R}^{N \times D}$ sebagai matriks data awal ($N$ baris data, masing-masing memiliki $D$ kolom fitur).

Operasi tingkat \textit{batch}:
\begin{align}
    \underbrace{Z^{(\ell)}}_{(N \times n_\ell)} &= \underbrace{X^{(\ell-1)}}_{(N \times n_{\ell-1})} \times \underbrace{W^{(\ell)}}_{(n_{\ell-1} \times n_\ell)} + \underbrace{\mathbf{1}_N}_{(N \times 1)} \times \underbrace{\mathbf{b}^{(\ell)}}_{(1 \times n_\ell)} \label{eq:batch_row_z} \\
    \underbrace{X^{(\ell)}}_{(N \times n_\ell)} &= \sigma^{(\ell)}\left( Z^{(\ell)} \right) \label{eq:batch_row_x}
\end{align}
dengan $\mathbf{1}_N = [1, 1, \dots, 1]^T \in \mathbb{R}^{N \times 1}$ adalah vektor kolom bernilai satu. Baris ke-$i$ dari matriks $Z^{(\ell)}$ dan $X^{(\ell)}$ masing-masing menyatakan $\mathbf{z}^{(\ell)}_i \in \mathbb{R}^{1 \times n_\ell}$ dan $\mathbf{x}^{(\ell)}_i \in \mathbb{R}^{1 \times n_\ell}$ dari data ke-$i$.

% -------------------------------------------------------------------------
\subsection{Fungsi Objektif}
\label{subsec:fungsi_objektif}

Pada layer terakhir ($\ell = L$), tidak diterapkan fungsi aktivasi non-linier, sehingga prediksi output berpadanan langsung dengan transformasi matriks:
\begin{equation}
    \hat{Y} = X^{(L-1)} W^{(L)} + \mathbf{1}_N \mathbf{b}^{(L)} \in \mathbb{R}^{N \times n_L} \label{eq:output_pred_batch}
\end{equation}
dan untuk data tunggal ke-$i$:
\begin{equation}
    \hat{\mathbf{y}}_i = \mathbf{x}_i^{(L-1)} W^{(L)} + \mathbf{b}^{(L)} \in \mathbb{R}^{1 \times n_L} \label{eq:output_pred_single}
\end{equation}
Diberikan matriks target $Y \in \mathbb{R}^{N \times n_L}$. Fungsi objektif \textit{Mean Squared Error} (MSE) dievaluasi melalui selisih target $Y$ dengan prediksi $\hat{Y}$:
\begin{equation}
    \mathcal{L} = \frac{1}{2N} \left\| Y - \hat{Y} \right\|_F^2 = \frac{1}{2N} \sum_{i=1}^N \sum_{j=1}^{n_L} \left( y_{i, j} - \hat{y}_{i, j} \right)^2 \label{eq:loss_mse}
\end{equation}
di mana $\|\cdot\|_F$ menyatakan norma Frobenius matriks.

% =========================================================================
% 2. PENURUNAN GRADIEN PARAMETER
% =========================================================================
\clearpage
\section{Backward Propagation}
\label{sec:penurunan_gradien}
\enlargethispage{2\baselineskip}

\subsection{Motivasi}
\label{subsec:motivasi_gd}

Pembaruan parameter matriks bobot $W^{(\ell)} \in \mathbb{R}^{n_{\ell-1} \times n_\ell}$ dan vektor bias $\mathbf{b}^{(\ell)} \in \mathbb{R}^{1 \times n_\ell}$ pada algoritma \textit{Gradient Descent} dengan \textit{learning rate} $\eta > 0$ dirumuskan sebagai:
\begin{align}
    W^{(\ell)}_{n+1} &= W^{(\ell)}_n - \eta \frac{\partial \mathcal{L}}{\partial W^{(\ell)}_n} \label{eq:update_w} \\
    \mathbf{b}^{(\ell)}_{n+1} &= \mathbf{b}^{(\ell)}_n - \eta \frac{\partial \mathcal{L}}{\partial \mathbf{b}^{(\ell)}_n} \label{eq:update_b}
\end{align}
dengan $n$ menyatakan indeks iterasi.

\subsection{Turunan Parameter pada Layer Output (\texorpdfstring{$\ell = L$}{l = L})}
\label{subsec:gradien_layer_L}

Karena pada layer output tidak diterapkan fungsi aktivasi, prediksi dinyatakan sebagai:
\begin{equation}
    \hat{Y} = X^{(L-1)} W^{(L)} + \mathbf{1}_N \mathbf{b}^{(L)} \in \mathbb{R}^{N \times n_L} \label{eq:pred_output_L}
\end{equation}

\begin{claim}
Turunan fungsi objektif $\mathcal{L}$ terhadap matriks prediksi $\hat{Y} \in \mathbb{R}^{N \times n_L}$ adalah:
\begin{equation}
    \frac{\partial \mathcal{L}}{\partial \hat{Y}} = \frac{1}{N} (\hat{Y} - Y) \in \mathbb{R}^{N \times n_L} \label{eq:claim_dL_dYhat}
\end{equation}
\end{claim}
\begin{proof}
Turunan fungsi skalar $\mathcal{L}$ terhadap matriks $\hat{Y}$ didefinisikan sebagai:
\begin{equation}
    \frac{\partial \mathcal{L}}{\partial \hat{Y}} = \begin{bmatrix}
        \frac{\partial \mathcal{L}}{\partial \hat{\mathbf{y}}_1} \\
        \frac{\partial \mathcal{L}}{\partial \hat{\mathbf{y}}_2} \\
        \vdots \\
        \frac{\partial \mathcal{L}}{\partial \hat{\mathbf{y}}_i} \\
        \vdots \\
        \frac{\partial \mathcal{L}}{\partial \hat{\mathbf{y}}_N}
    \end{bmatrix} = \begin{bmatrix}
        \frac{\partial \mathcal{L}}{\partial \hat{y}_{1, 1}} & \frac{\partial \mathcal{L}}{\partial \hat{y}_{1, 2}} & \dots & \frac{\partial \mathcal{L}}{\partial \hat{y}_{1, n_L}} \\
        \frac{\partial \mathcal{L}}{\partial \hat{y}_{2, 1}} & \frac{\partial \mathcal{L}}{\partial \hat{y}_{2, 2}} & \dots & \frac{\partial \mathcal{L}}{\partial \hat{y}_{2, n_L}} \\
        \vdots & \vdots & \ddots & \vdots \\
        \frac{\partial \mathcal{L}}{\partial \hat{y}_{i, 1}} & \frac{\partial \mathcal{L}}{\partial \hat{y}_{i, 2}} & \dots & \frac{\partial \mathcal{L}}{\partial \hat{y}_{i, n_L}} \\
        \vdots & \vdots & \ddots & \vdots \\
        \frac{\partial \mathcal{L}}{\partial \hat{y}_{N, 1}} & \frac{\partial \mathcal{L}}{\partial \hat{y}_{N, 2}} & \dots & \frac{\partial \mathcal{L}}{\partial \hat{y}_{N, n_L}}
    \end{bmatrix} \in \mathbb{R}^{N \times n_L}
\end{equation}
Dengan \textit{chain rule}, turunan terhadap setiap elemen $\hat{y}_{i, j}$:
\begin{equation}
    \frac{\partial \mathcal{L}}{\partial \hat{y}_{i, j}} = \frac{\partial \mathcal{L}}{\partial (y_{i, j} - \hat{y}_{i, j})} \frac{\partial (y_{i, j} - \hat{y}_{i, j})}{\partial \hat{y}_{i, j}} = \frac{1}{N} (y_{i, j} - \hat{y}_{i, j}) \cdot (-1) = \frac{1}{N} (\hat{y}_{i, j} - y_{i, j})
\end{equation}
Sehingga diperoleh $\frac{\partial \mathcal{L}}{\partial \hat{Y}} = \frac{1}{N} (\hat{Y} - Y)$.
\end{proof}

\begin{claim}
Gradien fungsi objektif terhadap matriks bobot $W^{(L)}$ adalah:
\begin{equation}
    \frac{\partial \mathcal{L}}{\partial W^{(L)}} = \frac{1}{N} \left( X^{(L-1)} \right)^T \left( \hat{Y} - Y \right) \in \mathbb{R}^{n_{L-1} \times n_L} \label{eq:claim_grad_W_L}
\end{equation}
\end{claim}
\begin{proof}
Dengan \textit{chain rule} terhadap $\hat{Y} = X^{(L-1)} W^{(L)} + \mathbf{1}_N \mathbf{b}^{(L)}$:
\begin{equation}
    \frac{\partial \mathcal{L}}{\partial W^{(L)}} = \frac{\partial \mathcal{L}}{\partial \hat{Y}} \frac{\partial \hat{Y}}{\partial W^{(L)}} = \left( X^{(L-1)} \right)^T \frac{\partial \mathcal{L}}{\partial \hat{Y}} = \frac{1}{N} \left( X^{(L-1)} \right)^T \left( \hat{Y} - Y \right) \qedhere
\end{equation}
\end{proof}

\begin{claim}
Gradien fungsi objektif terhadap vektor bias $\mathbf{b}^{(L)}$ adalah:
\begin{equation}
    \frac{\partial \mathcal{L}}{\partial \mathbf{b}^{(L)}} = \frac{1}{N} \mathbf{1}_N^T \left( \hat{Y} - Y \right) \in \mathbb{R}^{1 \times n_L} \label{eq:claim_grad_b_L}
\end{equation}
\end{claim}
\begin{proof}
Dengan \textit{chain rule} terhadap $\hat{Y} = X^{(L-1)} W^{(L)} + \mathbf{1}_N \mathbf{b}^{(L)}$:
\begin{equation}
    \frac{\partial \mathcal{L}}{\partial \mathbf{b}^{(L)}} = \frac{\partial \mathcal{L}}{\partial \hat{Y}} \frac{\partial \hat{Y}}{\partial \mathbf{b}^{(L)}} = \mathbf{1}_N^T \frac{\partial \mathcal{L}}{\partial \hat{Y}} = \frac{1}{N} \mathbf{1}_N^T \left( \hat{Y} - Y \right) \qedhere
\end{equation}
\end{proof}

\clearpage
\subsection{Turunan Parameter pada Layer \texorpdfstring{$\ell = L-1, \dots, 2, 1$}{l = L-1, ..., 2, 1}}
\label{subsec:gradien_layer_hidden}

Pada hidden layer $\ell \in \{L-1, \dots, 2, 1\}$, fungsi aktivasi $\sigma^{(\ell)}$ diterapkan terhadap bentuk matriks:
\begin{equation}
    X^{(\ell)} = \sigma^{(\ell)}\left( X^{(\ell-1)} W^{(\ell)} + \mathbf{1}_N \mathbf{b}^{(\ell)} \right) \in \mathbb{R}^{N \times n_\ell} \label{eq:X_hidden_matrix}
\end{equation}

\noindent \textbf{\textit{Backward Propagation} pada Layer $\ell = L-1$}: \\
\textit{Backward propagation} galat dari layer output ke fitur $X^{(L-1)}$ dievaluasi melalui aturan rantai:
\begin{equation}
    \frac{\partial \mathcal{L}}{\partial X^{(L-1)}} = \frac{\partial \mathcal{L}}{\partial \hat{Y}} \frac{\partial \hat{Y}}{\partial X^{(L-1)}} = \frac{1}{N} \left( \hat{Y} - Y \right) \left( W^{(L)} \right)^T \in \mathbb{R}^{N \times n_{L-1}} \label{eq:dL_dXLprev}
\end{equation}
Turunan terhadap bentuk pre-aktivasi layer $L-1$:
\begin{equation}
    \frac{\partial \mathcal{L}}{\partial \left( X^{(L-2)} W^{(L-1)} + \mathbf{1}_N \mathbf{b}^{(L-1)} \right)} = \frac{\partial \mathcal{L}}{\partial X^{(L-1)}} \odot {\sigma^{(L-1)}}'\left( X^{(L-2)} W^{(L-1)} + \mathbf{1}_N \mathbf{b}^{(L-1)} \right)
\end{equation}
Gradien parameter $W^{(L-1)}$ dan $\mathbf{b}^{(L-1)}$:
\begin{align}
    \frac{\partial \mathcal{L}}{\partial W^{(L-1)}} &= \left( X^{(L-2)} \right)^T \frac{\partial \mathcal{L}}{\partial \left( X^{(L-2)} W^{(L-1)} + \mathbf{1}_N \mathbf{b}^{(L-1)} \right)} \in \mathbb{R}^{n_{L-2} \times n_{L-1}} \\[0.3em]
    \frac{\partial \mathcal{L}}{\partial \mathbf{b}^{(L-1)}} &= \mathbf{1}_N^T \frac{\partial \mathcal{L}}{\partial \left( X^{(L-2)} W^{(L-1)} + \mathbf{1}_N \mathbf{b}^{(L-1)} \right)} \in \mathbb{R}^{1 \times n_{L-1}}
\end{align}
\textit{Backward propagation} ke fitur layer $L-2$:
\begin{equation}
    \frac{\partial \mathcal{L}}{\partial X^{(L-2)}} = \frac{\partial \mathcal{L}}{\partial \left( X^{(L-2)} W^{(L-1)} + \mathbf{1}_N \mathbf{b}^{(L-1)} \right)} \left( W^{(L-1)} \right)^T \in \mathbb{R}^{N \times n_{L-2}}
\end{equation}

\noindent \textbf{\textit{Backward Propagation} pada Layer $\ell = L-2$}: \\
Turunan terhadap bentuk pre-aktivasi layer $L-2$:
\begin{equation}
    \frac{\partial \mathcal{L}}{\partial \left( X^{(L-3)} W^{(L-2)} + \mathbf{1}_N \mathbf{b}^{(L-2)} \right)} = \frac{\partial \mathcal{L}}{\partial X^{(L-2)}} \odot {\sigma^{(L-2)}}'\left( X^{(L-3)} W^{(L-2)} + \mathbf{1}_N \mathbf{b}^{(L-2)} \right)
\end{equation}
Gradien parameter $W^{(L-2)}$ dan $\mathbf{b}^{(L-2)}$:
\begin{align}
    \frac{\partial \mathcal{L}}{\partial W^{(L-2)}} &= \left( X^{(L-3)} \right)^T \frac{\partial \mathcal{L}}{\partial \left( X^{(L-3)} W^{(L-2)} + \mathbf{1}_N \mathbf{b}^{(L-2)} \right)} \in \mathbb{R}^{n_{L-3} \times n_{L-2}} \\[0.3em]
    \frac{\partial \mathcal{L}}{\partial \mathbf{b}^{(L-2)}} &= \mathbf{1}_N^T \frac{\partial \mathcal{L}}{\partial \left( X^{(L-3)} W^{(L-2)} + \mathbf{1}_N \mathbf{b}^{(L-2)} \right)} \in \mathbb{R}^{1 \times n_{L-2}}
\end{align}
\textit{Backward propagation} ke fitur layer $L-3$:
\begin{equation}
    \frac{\partial \mathcal{L}}{\partial X^{(L-3)}} = \frac{\partial \mathcal{L}}{\partial \left( X^{(L-3)} W^{(L-2)} + \mathbf{1}_N \mathbf{b}^{(L-2)} \right)} \left( W^{(L-2)} \right)^T \in \mathbb{R}^{N \times n_{L-3}}
\end{equation}

\noindent \textbf{\textit{Backward Propagation} Umum pada Hidden Layer ($\ell = L-1, \dots, 2, 1$)}: \\
Secara umum, untuk setiap hidden layer $\ell \in \{L-1, \dots, 2, 1\}$:
\begin{align}
    \frac{\partial \mathcal{L}}{\partial \left( X^{(\ell-1)} W^{(\ell)} + \mathbf{1}_N \mathbf{b}^{(\ell)} \right)} &= \frac{\partial \mathcal{L}}{\partial X^{(\ell)}} \odot {\sigma^{(\ell)}}'\left( X^{(\ell-1)} W^{(\ell)} + \mathbf{1}_N \mathbf{b}^{(\ell)} \right) \in \mathbb{R}^{N \times n_\ell} \label{eq:batch_dL_hidden} \\[0.3em]
    \frac{\partial \mathcal{L}}{\partial W^{(\ell)}} &= \left( X^{(\ell-1)} \right)^T \frac{\partial \mathcal{L}}{\partial \left( X^{(\ell-1)} W^{(\ell)} + \mathbf{1}_N \mathbf{b}^{(\ell)} \right)} \in \mathbb{R}^{n_{\ell-1} \times n_\ell} \label{eq:batch_grad_w_hidden} \\[0.3em]
    \frac{\partial \mathcal{L}}{\partial \mathbf{b}^{(\ell)}} &= \mathbf{1}_N^T \frac{\partial \mathcal{L}}{\partial \left( X^{(\ell-1)} W^{(\ell)} + \mathbf{1}_N \mathbf{b}^{(\ell)} \right)} \in \mathbb{R}^{1 \times n_\ell} \label{eq:batch_grad_b_hidden} \\[0.3em]
    \frac{\partial \mathcal{L}}{\partial X^{(\ell-1)}} &= \frac{\partial \mathcal{L}}{\partial \left( X^{(\ell-1)} W^{(\ell)} + \mathbf{1}_N \mathbf{b}^{(\ell)} \right)} \left( W^{(\ell)} \right)^T \in \mathbb{R}^{N \times n_{\ell-1}} \label{eq:batch_dL_dXprev}
\end{align}
dengan kondisi batas input awal $X^{(0)} = X \in \mathbb{R}^{N \times D}$.

\end{document}
